\chapter*{Resumen}
\addcontentsline{toc}{chapter}{Resumen}

Angine de Poitrine es un dúo de Saguenay (Quebec) que toca con guitarra y bajo de trastes de cuarto de tono, construye
sus canciones con un looper y usa compases poco comunes. Este trabajo pregunta si su música encaja en un estilo
existente o si es una propuesta nueva, y responde la pregunta midiendo desde el audio. Para ello se construyó una
herramienta abierta que transcribe en 24-EDO (veinticuatro divisiones iguales de la octava): separa la guitarra y el
bajo, estima la altura con una saliencia armónica en una rejilla de 25 cents, usa la repetición del looper para
estabilizar la transcripción, segmenta la canción en compases y asigna a cada compás un acorde con las calidades de
Piston extendidas a 24-EDO, sus funciones (Riemann) y sus ejes (Lendvai). La salida principal es el \emph{tabroll}, un
piano roll de 24 filas por octava con los acordes debajo.

La herramienta se evaluó contra partituras corregidas de oído, con arbitrajes a ciegas y una canción reservada como
prueba ciega. La fundamental del acorde acierta en promedio el 85\,\% en los compases confiables, mientras que la
melodía acierta el 49\,\%; la diferencia se explica por la forma en que se decide la armonía (peso del bajo, suma
sobre las vueltas del loop, independencia de la octava y resolución por compás). Con esa base, cada canción del álbum
\emph{Vol.~II} se resumió en cuatro bloques de rasgos (armonía, escala, cuartos de tono y metro) y se comparó con seis
corpus de estilos mediante la distancia de Jensen-Shannon y un percentil de tipicidad, después de aplicar a los corpus
el error de medición de la propia herramienta, estimado con canciones de control.

Los resultados indican que la banda es un híbrido con un elemento propio. En armonía, ninguna canción es típica de
ningún corpus (máximo 17.7 sobre 50) y el 37.5\,\% de los casos queda sin vecino, frente al 5.2\,\% en piezas reales de
rock. En melodía, dos grados en cuarto de tono concentran entre el 48 y el 82\,\% de los cuartos de tono de cada
canción: son grados fijos de la escala, no adornos, y forman dos familias (frigia y dórica), de las cuales la dórica no
aparece en ningún corpus medido. En metro, los ciclos del looper muestran periodos de 20 y 28 corcheas, compatibles con
compases de 5 y de 7 tiempos. La conclusión es que la música de la banda no corresponde a un estilo existente, sino a
una recombinación de lenguajes conocidos con un rasgo nuevo: los cuartos de tono como grados fijos de la escala.

\medskip
\noindent\textbf{Palabras clave:} microtonalidad, 24-EDO, transcripción automática, recuperación de información musical,
armonía funcional, corpus de estilos, divergencia de Jensen-Shannon, Angine de Poitrine.

\chapter*{Abstract}
\addcontentsline{toc}{chapter}{Abstract}

\pendiente{traducir el resumen al inglés cuando el texto en español quede estable.}
